\addchap{Vorwort}
\begin{refsection}

Lehrkräfte an Schulen haben – unabhängig vom Fach, das sie unterrichten – ständig mit Sprache zu tun. Zu den Fragen, die sich dabei stellen, gehören: Warum kann man im Deutschen nicht sagen: \textit{Ich lächle dich}, obwohl das im Französischen funktioniert? Warum sagt man im bundesdeutschen Standard nicht \textit{Heute ich habe gegessen einen Apfel}, obwohl das im Englischen funktioniert? Warum hört man von Kindern Äußerungen wie \textit{mit dem Affe} statt \textit{mit dem Affen}? Warum gibt es vergleichsweise viele Kinder, die das [s] zeitweise schwierig zu produzieren finden, obwohl sie andere Laute perfekt sprechen können? Und warum haben einige Menschen das Gefühl, sie sollten das Wort \textit{einige}, wenn kein Substantiv folgt, vielleicht groß schreiben: \textit{Dieses Gefühl haben Einige}? 

Damit Lehrkräfte auf solche Fragen vorbereitet sind, ist Sprachbildung inzwischen fest im Lehramtsstudium verankert. Das vorliegende Buch unterstützt die Ausbildung in diesem Bereich und bietet eine fundierte Einführung in die sprachlichen Einheiten und Strukturen des Deutschen. Zentrale Konzepte und Begriffe werden in hervorgehobenen Boxen erläutert. Die Beispiele und Analysen basieren überwiegend auf authentischen Belegen aus Texten der bundesdeutschen Standardschriftlichkeit sowie dem gesprochenen Standard. Ergänzend werden auch die sprachliche Variation und die Entwicklung von Strukturen thematisiert.     

Das Buch richtet sich in erster Linie an Studierende des Lehramts Deutsch für Grund- und Oberschulen. In jedem Kapitel wird die schulische Relevanz der behandelten Themen explizit aufgezeigt. 

In erster Linie ist das Buch als vertiefendes Begleitmaterial zu einführenden Lehrveranstaltungen konzipiert. Die Idee für dieses Lehrbuch entstand aus zahlreichen lehramtsspezifischen Veranstaltungen an der Humboldt-Universität zu Berlin, bei denen aufgrund zeitlicher Einschränkungen oft nicht genügend Raum für weiterführende Diskussionen oder die praktische Anwendung von Analysetechniken blieb. Aus diesem Grund werden zu jedem Thema verschiedene Analysetechniken ausführlich vorgestellt und begründet. In zahlreichen gesonderten Exkursen werden einzelne Themen vertieft, alternative Erklärungen angeboten und schulrelevante Aspekte diskutiert. So kann das Buch auch in späteren Seminaren zur Vertiefung genutzt werden.    

Das Lehrbuch eignet sich zudem für das Selbststudium, etwa im Rahmen einer Abschlussarbeit oder als Weiterbildungsressource für Lehrkräfte. 

Bei lehrveranstaltungsbegleitender Nutzung empfehlen wir ein kapitelweises Lesen. Beispielanalysen eröffnen die Möglichkeit, das Gelesene praktisch anzuwenden und zu üben. Die Inhalte eines jeden Kapitels werden am Ende zusammengefasst. Abschließend bieten jeweils vier Aufgaben die Möglichkeit, das erworbene Wissen anzuwenden und durch den Vergleich mit Musterlösungen übend zu lernen. Nutzerinnen und Nutzer, die sich für einzelne Aspekte interessieren, das Buch also zum Nachschlagen benutzen, können das Thema, das sie interessiert, sowohl im Inhaltsverzeichnis als auch im ausführlichen Register finden. Einzelne Abschnitte und ganze Kapitel sind durch entsprechende Verweise so konzipiert, dass sie auch unabhängig voneinander verständlich sind. In der pdf-Version sind jene Verweise verlinkt, sodass ein Klick genügt, um den relevanten Textabschnitt zu finden. Zahlreiche Passagen sind als Exkurse ausgezeichnet. Die Exkurse vertiefen Aspekte entweder inhaltlich oder sie stellen einen direkten Schulbezug her und können übersprungen werden, ohne dass das Verständnis der folgenden Abschnitte gefährdet ist.

Wer nach einer Einführung in die Linguistik sucht, wird eine Fülle an Lehrwerken finden, auch viele deutlich weniger umfangreiche. Das liegt einerseits daran, dass wir einen recht breiten Adressatenkreis ansprechen. Andererseits sind die Erwartungen an Lehrkräfte aufgrund der sprachlich immer heterogeneren Schülerschaft gestiegen. Lehrkräfte benötigen ein weit tieferes Wissen, als das, welches sie den Schülerinnen und Schülern vermitteln. Die Inhalte des Buches gehen deshalb weit über die Schulgrammatik hinaus.

Unser Erklärungsansatz ist grammatisch-funktional: grammatisch, weil wir bei der Form sprachlicher Einheiten ansetzen – und funktional, weil wir nach der sprachlichen Funktion jener Einheiten innerhalb von größeren sprachlichen Einheiten fragen. Deshalb bewegen wir uns von den größeren linguistischen Einheiten (den Sätzen) zu den kleineren (Wortgruppen, Wörtern, Silben, Lauten). Das abschließende Kapitel zum Schriftsystem des Deutschen (Graphematik) greift verschiedene Aspekte der vorangegangenen Kapitel auf und erläutert deren Relevanz für den Schriftspracherwerb.

Bei der Auswahl der Termini orientieren wir uns weitgehend an dem \citet{KMK2019}, das die Kultusministerkonferenz 2019 zustimmend zur Kenntnis genommen hat. Dieses Verzeichnis bildet die wesentliche fachwissenschaftliche Grundlage für den Grammatikunterricht und die grammatische Sprachreflexion in der Schule, und zwar unabhängig vom Schultyp.

Immer wieder weisen wir auf Schnittstellen zum schulischen Grammatik- und Sprachunterricht, zur schulischen Reflexion über Grammatik hin. Aber es handelt sich hier weder um ein Kompendium der zu unterrichtenden Inhalte noch um didaktische Konzepte ihrer schulischen Vermittlung. 

Jede grammatische Beschreibung bedarf theoretischer Beschreibungsmodelle. Allerdings stehen entsprechende Theorien hier nicht im Vordergrund. Sie werden kurz vorgestellt und, wo nötig, diskutiert. Durchgängig bieten detaillierte Literaturverweise, einige Exkurse und Weiterführendes in Fußnoten Zugang zu einer tieferen Beschäftigung mit einzelnen Phänomenen und ihrem Wandel, zum Beispiel bei einer sprachwissenschaftlichen Bachelor- oder Masterarbeit.


Sprache und Sprechen sind viel mehr als das hier Dargestellte. Der lebendige und angemessene, schöne oder als weniger schön geltende, alltägliche oder poetische Gebrauch der sich permanent wandelnden Sprache ist vielfältig, spielerisch und basiert auf viel mehr und teilweise ganz anderen als analytischen grammatischen Fähigkeiten, die hier im Zentrum stehen. Wo möglich, weisen wir auf die Vielfalt sprachlicher Ausdrucksmöglichkeiten hin. Die Strukturbeschreibungen setzen mehrheitlich bei authentischen Sprachdaten der geschriebenen Sprache aus Texten des sogenannten Standards an. 

Wir wollen mit unserem Buch alle Geschlechter gleichermaßen ansprechen. Da sich diesbezüglich bis heute kein Standard etabliert hat und jede Form des Genderns sowohl Zuspruch als auch Ablehnung erfährt, verwenden wir unterschiedliche Arten der Nennung, etwa Partizipien, Beidnennung, Klammerlösungen und Sachbezeichnungen. Unser Ziel dabei ist, den Text angenehm lesbar zu gestalten.


Wir danken insbesondere den Studierenden für ihre Rückmeldungen und die anregenden Diskussionen, die dieses Lehrbuch seit seiner Entstehung als Manuskript begleitet und geprägt haben. Unser Dank gilt unseren Kollegen des Instituts für deutsche Sprache und Linguistik für Anregungen und Diskussionen, insbesondere Antje Sauermann, Eva Schlachter, Malte Belz, Antonio Machicao y Priemer, Hagen Hirschmann, Stefan Müller, Christine Mooshammer, Klaus Welke, Anke Lüdeling und Sarah Zobel. Gedankt ist den Herausgebern der Reihe Textbooks in Language Sciences Antonio Machicao y Priemer und Stefan Müller. Für das Korrekturlesen danken wir herzlich Franziska Dilling und Torsten Föste. Nicht zuletzt sei den anonymen Gutachterinnen und Gutachtern für Ideen und Korrekturen gedankt. 
\\
\\
\\


\noindent Berlin, am \today \hfill Marc Felfe und Jana Brunner



%\printbibliography[heading=subbibliography]
\end{refsection}